% !PRINTED 316
\weberpage{316}

Man kann ihr aber noch eine andere, sehr bemerkenswerte Form geben.

Nach den zwischen den Funktionen $f(\omega), f_1(\omega), f_2(\omega)$
bestehenden Relationen [§\,34, (11), (12)] ist:

\[
\frac{f(\omega)^{24} - 64}{f(\omega)^{12}} = \frac{[f_1(\omega)^8 - f_2(\omega)^8]^2}{f(\omega)^4}.
\]

Führen wir also für $w$ die neue Unbekannte

\weq{26}{y = \sqrt{w} = \frac{f_1(\omega)^8 - f_2(\omega)^8}{f(\omega)^2(w^2+5)}}

ein, so erhalten wir die Form:

\weq{27}{y^5 + 5y = \frac{f_1(\omega)^8 - f_2(\omega)^8}{f(\omega)^2}.}

Auf diese Formeln gründet sich die von \gesperrt{Hermite} und
\gesperrt{Kronecker} geschaffene Auflösung der Gleichung fünften Grades
durch elliptische Funktionen, auf die wir im 14. Abschnitt des II.
Bandes hingewiesen haben.

Nach Bd.\,I, §\,60 kann die allgemeine Gleichung fünften Grades auf die
Form reduziert werden:

\weq{28}{z^5 + 5z = a,}

und diese Gleichung wird mit (27) identisch, wenn man setzt:

\[
f_1(\omega)^8 - f_2(\omega)^8 = af^2(\omega).
\]

Nimmt man hierzu die Gleichungen:

\[
f_1^8 = f_2^8 = f^8, \quad ff_1f_2 = \sqrt{2}, % [sic] Die Quelle druckt beide Gleichheitszeichen.
\]

so erhält man

\weq{29}{f^{24} - a^2f^{12} - 64 = 0.}

Durch diese quadratische Gleichung ist $f^{12}$ als Funktion von $a$
bestimmt.

Dann sind die Wurzeln der Gleichung (28) durch die Quadratwurzeln aus
den Ausdrücken (22) dargestellt. Das doppelte Vorzeichen dieser
Quadratwurzeln erklärt sich daraus, daß zwei entgegengesetzte Werte von
$a$ zu derselben Gleichung (29) führen. Man kann aber auch die Wurzeln
eindeutig bestimmen nach der Formel (26):

\weq{30}{z = \frac{a}{w^2+5}.}
